\section{WOTS+C}
The standard Winternitz One-Time Signature (WOTS+) is the building block of most modern hash-based schemes. It works by encoding the message into a series of small integers and using them to iterate hash chains \cite{wots+}. A significant portion of a WOTS+ signature size (approx 20\%) is dedicated to "Checksum," which is necessary to prevent a forgery attack where an adversary extends the hash chains forward. WOTS+C (compressed WOTS+) eliminates the transmission of this checksum entirely, achieving a smaller signature size. 

In WOTS+, the checksum $C$ is calculated from the message digest $D$.

$$C = \sum_{i=1}^{l} (w - 1 - d_i)$$

This checksum is then signed using additional chains. In WOTS+C\cite{wots+c}, instead of calculating a checksum  for a fixed message digest, the signer modifies the message digest until it possesses a specific,  pre-determined checksum property. This is done by appending a counter \texttt{ctr} and a random salt \texttt{R} to the message before hashing. The signer increments \texttt{ctr} until the resulting hash digest $D$ satisfies a global target condition.

\subsection{The Target Condition}
WOTS utilize a parameter \texttt{w} to define the number of hashchains and correspondingly their length. The message digest length is \texttt{n} bytes so \texttt{l = 8n/(log2(w))}. The values $d_{i\in[l]}$ are treated as integers in \texttt{0...(w-1)}. 

Let \texttt{Swn} be a protocol parameter (often near the expected value, but may be chosen larger to trade signing/grinding cost for verification cost). If the hash output is uniform, the expected value of a single digit is \texttt{d\_i = (w-1)/2}. For \texttt{l} digits, the expected \texttt{Swn} is: 

$$S_{wn} = \sum_{i=0}^{l-1}d_i$$

By enforcing this condition during signing (including grinding), the verifier can simply assume the checksum is \texttt{Swn} without receiving it. If the reconstructed digest does not sum to \texttt{Swn}, the signature is invalid. 

Additionally we can manipulate a target condition to decrease the cost of signature verification. The signature generation and verification operations together consist of \texttt{St = (w-1)*l} hashing operations. So we can increase \texttt{Swn} (grinding complexity) and receive the reduced verification complexity of \texttt{(St - Swn)} hashes. Concrete parameters are presented in Section~4.

\subsection{Base-w Representation}
WOTS represents a message digest in base \texttt{w} so that each digit is an integer in \texttt{[0...w-1]}. These digits directly determine the number of hashes to apply to each WOTS chain element during signing and verification. The \texttt{base\_w()} function converts a byte string into an array of \texttt{out\_len} base-w digits (this requires \texttt{w} to be a power of two.).
\begin{verbatim}
    Algorithm: base_w(m, w, out_len) 
        m: byte string
        w: Winternitz parameter
        out_len: output length
    Output: 
        b_w: array of out_len integers in [0, w-1]
    
    1. in_idx ← 0
    2. bits ← 0
    3. total ← 0
    4. b_w ← []
    5. for i from 0 to (out_len - 1):
        a. if bits == 0:
            total ← m[in_idx]
            in_idx ← in_idx + 1
            bits ← 8
        b. bits ← bits - log2(w)
        c. b_w[i] ← (total >> bits) AND (w - 1)
    6. return b_w
\end{verbatim} 

We provide optimized versions for \texttt{w = 256} and \texttt{w = 4} proposed by different SHRINCS instances. For \texttt{w=256}:
\begin{verbatim}
    Algorithm: base_w256(m, out_len)
    Input:
        m: byte string
        out_len: output length 
    Output:
        b_w: array of out_len integers in [0, 255]

    1. b_w ← []
    2. for i from 0 to (out_len - 1):
        a. b_w[i] ← m[i]
    3. return b_w
\end{verbatim}
For \texttt{w=4}:
\begin{verbatim}
    Algorithm: base_w4(m, out_len)
    Input:
        m: byte string
        out_len: output length
    Output:
        b_w: array of out_len integers in [0,3]

    1. b_w ← []
    2. for j from 0 to (out_len/4 - 1):
        a. x ← m[j]
        b. b_w[4*j + 0] ← (x >> 6) AND 3
        c. b_w[4*j + 1] ← (x >> 4) AND 3
        d. b_w[4*j + 2] ← (x >> 2) AND 3
        e. b_w[4*j + 3] ← x AND 3
    3. return b_w
\end{verbatim}

\subsection{Chain Function}
WOTS derives signature components and public key material by iteratively applying the tweakable hash function along a Winternitz chain. The function \texttt{chain()} takes an n-byte input value \texttt{m} and applies \texttt{Tw(PK.seed, ADRS, ·)} exactly \texttt{steps} times, beginning at chain position \texttt{start}. For each iteration, the address field hash inside \texttt{ADRS} is set to the current chain position \texttt{i} (from \texttt{start} to \texttt{start+steps-1}) to ensure domain separation between different chain steps. The output is the resulting n-byte value after advancing the chain by \texttt{steps}, which is used both for producing signature chain values (advancing from 0 to \texttt{msg[i]}) and for reconstructing end-of-chain values during verification (advancing from \texttt{msg[i]} to \texttt{w-1}).
\begin{verbatim}
    Algorithm: chain(m, start, steps, PK.seed, ADRS)
    Input:
        m: input
        start: starting index
        steps: number of iterations
        PK.seed: public seed
        ADRS: tweak
    Output:
        n-byte chain output

    1. tmp ← m
    2. for i from start to (start + steps - 1):
        a. setHashAddress(ADRS, i)
        b. tmp ← Tw(PK.seed, ADRS, tmp)
    3. return tmp
\end{verbatim}

\subsection{Key Generation}
To create a WOTS public key for a given \texttt{keypair} index, we derive one secret starting value per chain using \texttt{PRF}, then advance each value to the end of its chain. The resulting array of end-of-chain values is compressed to a fixed-size WOTS public key for that \texttt{keypair}.
\begin{verbatim}
    Algorithm: wots_PKgen(SK.seed, PK.seed, ADRS, keypair)
    Input:
        SK.seed: secret seed
        PK.seed: public seed
        ADRS: tweak
        keypair: keypair index
    Output:
        pk: compressed WOTS+C public key

    1. tmp ← []
    2. for i from 0 to l - 1:
        a. setTypeAndClear(ADRS, PRF)
        b. setKeyPairAddress(ADRS, keypair)
        c. setChainAddress(ADRS, i)
        d. sk_i ← PRF(PK.seed, ADRS, SK.seed)
        e. setTypeAndClear(ADRS, SF_WOTS_CHAIN)
        f. setKeyPairAddress(ADRS, keypair)
        g. setChainAddress(ADRS, i)
        h. tmp[i] ← chain(sk_i, 0, w - 1, PK.seed, ADRS)
    3. setTypeAndClear(ADRS, SF_TREE)
    4. setKeyPairAddress(ADRS, keypair)
    5. pk ← Tw(PK.seed, ADRS, tmp[0] || ... || tmp[l-1])
    6. return pk
\end{verbatim}

\subsection{Messsage Digest Grinding}
WOTS+C enforces the checksum/target condition by searching over the counter \texttt{ctr}. The value \texttt{R} (a randomness/nonce input to message hashing) may be fixed (deterministic) and need not be resampled during grinding. The function \texttt{wots\_grind()} increments \texttt{ctr} and recomputes the message digest until the derived base-w digits satisfy the target condition. It outputs the found \texttt{ctr} and the corresponding digest (and \texttt{R} is carried in the signature only if verification requires it to recompute the same digest).

\oleks{To consider if we need to mention counter in Section 5.4}
\begin{verbatim}
    Algorithm: wots_grind(m, R, PK.seed, PK.root, ADRS)
    Input:
        m: message
        R: randomness
        PK.seed: public seed
        PK.root: public key root
        ADRS: tweak
    Output:
        msg: array of l integers
        ctr: final counter value used

    1. setTypeAndClear(ADRS, H_MSG)
    2. for ctr from 0 to 2^32 - 1:
        a. digest ← H_msg(ADRS, R || ctr, PK.seed, PK.root, m)
        b. msg ← base_w(digest, l)
        c. sum ← 0
        d. for i from 0 to (l - 1):
            sum ← sum + msg[i]
        e. if sum == Swn:
            return (msg, ctr)
    3. abort with error  // R must be resampled
\end{verbatim}
\oleks{to add expected grinding complexity for different approaches}

Additionally we need to define a func which allows to calculate the digest without grinding (by providing the \texttt{ctr} value directly).

\begin{verbatim}
    Algorithm: wots_digest(m, R, ctr, PK.seed, PK.root, ADRS)
    Input:
        m: message
        R: randomness
        ctr: counter
        PK.seed: public seed
        PK.root: public key root
        ADRS: tweak
    Output:
        msg: array of l integers
        valid: boolean indicating if sum == Swn

    1. setTypeAndClear(ADRS, H_MSG)
    2. digest ← H_msg(ADRS, R || ctr, PK.seed, PK.root, m)
    3. msg ← base_w(digest, l)
    4. sum ← 0
    5. for i from 0 to l - 1:
        sum ← sum + msg[i]
    6. valid ← (sum == Swn)
    7. return (msg, valid)
\end{verbatim}

\subsection{Signature Generation}
To sign, WOTS+C instance computes the message digits and then reveals, for each chain, the chain value after exactly \texttt{msg[i]} steps (starting from the secret-derived chain \texttt{start}). These revealed values form the WOTS signature vector. The signature also carries the randomness 
\texttt{R} and grinding counter \texttt{ctr} needed for the verifier to reproduce the same message digits.
\begin{verbatim}
    Algorithm: wots_sign(m, SK.seed, SK.prf, PK.seed, PK.root, ADRS, keypair)
    Input:
        m: message
        SK.seed: secret seed
        SK.prf: PRF key for randomness
        PK.seed: public seed
        PK.root: public key root
        ADRS: tweak
        keypair: keypair index
    Output:
        sig: WOTS+C signature (R || ctr || chain_values)

    1. opt_rand ← random(n) or PK.seed  // for randomized or deterministic signing
    2. R ← PRF_msg(SK.prf, PK.seed, opt_rand, m)
    3. (msg, ctr) ← wots_grind(m, R, PK.seed, PK.root, ADRS)
    4. sig ← R || toByte(ctr, c)
    5. for i from 0 to (l - 1):
        a. setTypeAndClear(ADRS, PRF)
        b. setKeyPairAddress(ADRS, keypair)
        c. setChainAddress(ADRS, i)
        d. sk_i ← PRF(PK.seed, ADRS, SK.seed)
        e. setTypeAndClear(ADRS, SF_WOTS_CHAIN)
        f. setKeyPairAddress(ADRS, keypair)
        g. setChainAddress(ADRS, i)
        h. sig ← sig || chain(sk_i, 0, msg[i], PK.seed, ADRS)
    6. return sig
\end{verbatim}

\subsection{Signature Verification (Public Key Recovery)}
Verification reconstructs the WOTS+C public key implied by the signature. For each chain element, the verifier continues hashing forward from the signature’s provided intermediate value for \texttt{w-1-msg[i]} steps to reach the end-of-chain value. Compressing all end-of-chain values yields the recovered WOTS+C public key, which can then be checked against the expected value in the surrounding scheme.
\begin{verbatim}
    Algorithm: wots_pkFromSig(sig, m, PK.seed, PK.root, ADRS, keypair)
    Input:
        sig: signature
        m: signed message
        PK.seed: public seed
        PK.root: public key root
        ADRS: tweak
        keypair: keypair index
    Output:
        pk: recovered public key or $\bot$ if invalid

    1. R ← sig[0:R_SIZE]                 // maybe it makes sense to define sig_parse op
    2. ctr ← sig[R_SIZE:R_SIZE + c]      
    3. chains_start ← R_SIZE + c 
    4. (msg, valid) ← wots_digest(m, R, ctr, PK.seed, PK.root, ADRS)
    5. if not valid:
        return $\bot$
    6. tmp ← []
    7. for i from 0 to (l - 1):
        a. setTypeAndClear(ADRS, SF_WOTS_CHAIN)
        b. setKeyPairAddress(ADRS, keypair)
        c. setChainAddress(ADRS, i)
        d. sig_i ← sig[chains_start + i * n : chains_start + (i+1)*n]
        e. tmp[i] ← chain(sig_i, msg[i], w - 1 - msg[i], PK.seed, ADRS)
    8. setTypeAndClear(ADRS, SF_TREE)
    9. setKeyPairAddress(ADRS, keypair)
    10. pk ← Tw(PK.seed, ADRS, tmp[0] || ... || tmp[l-1])
    11. return pk
\end{verbatim}